
\documentclass{article}
\usepackage[table]{xcolor}
\usepackage[left=2cm, right=2cm]{geometry}

\title{Assignment 3}
\author{Christian Parker \\ CXP22004}

\begin{document}

\maketitle

\begin{enumerate}
    % Problem 1
    \item \ % Add a space to force a left alignment for the number.
    \begin{center}
        \begin{tabular}{|c|c|c|c|c|c|c|c|c|c|c|}
        \multicolumn{11}{c}{\textbf{Dilation}} \\
        \hline
        0 & 0 & 0 & 0 & 0 & 0 & 0 & 0 & 0 & 0 & 0 \\
        \hline
        0 & 0 & 0 & 0 & 0 & 0 & 0 & 0 & 0 & 0 & 0 \\
        \hline
        0 & \cellcolor{blue!20}1 & \cellcolor{blue!20}1 & \cellcolor{blue!20}1 & \cellcolor{blue!20}1 & \cellcolor{blue!20}1 & \cellcolor{blue!20}1 & \cellcolor{blue!20}1 & \cellcolor{blue!20}1 & 0 & 0 \\
        \hline
        0 & \cellcolor{blue!20}1 & \cellcolor{blue!20}1 & \cellcolor{blue!20}1 & \cellcolor{blue!20}1 & \cellcolor{blue!20}1 & \cellcolor{blue!20}1 & \cellcolor{blue!20}1 & \cellcolor{blue!20}1 & 0 & 0 \\
        \hline
        0 & \cellcolor{blue!20}1 & \cellcolor{blue!20}1 & \cellcolor{blue!20}1 & \cellcolor{blue!20}1 & \cellcolor{blue!20}1 & \cellcolor{blue!20}1 & \cellcolor{blue!20}1 & \cellcolor{blue!20}1 & 0 & 0 \\
        \hline
        0 & \cellcolor{blue!20}1 & \cellcolor{blue!20}1 & \cellcolor{blue!20}1 & \cellcolor{blue!20}1 & \cellcolor{blue!20}1 & \cellcolor{blue!20}1 & \cellcolor{blue!20}1 & \cellcolor{blue!20}1 & 0 & 0 \\
        \hline
        0 & \cellcolor{blue!20}1 & \cellcolor{blue!20}1 & \cellcolor{blue!20}1 & 0 & \cellcolor{blue!20}1 & \cellcolor{blue!20}1 & \cellcolor{blue!20}1 & \cellcolor{blue!20}1 & 0 & 0 \\
        \hline
        0 & \cellcolor{blue!20}1 & \cellcolor{blue!20}1 & \cellcolor{blue!20}1 & \cellcolor{blue!20}1 & \cellcolor{blue!20}1 & \cellcolor{blue!20}1 & \cellcolor{blue!20}1 & \cellcolor{blue!20}1 & 0 & 0 \\
        \hline
        0 & 0 & \cellcolor{blue!20}1 & \cellcolor{blue!20}1 & \cellcolor{blue!20}1 & \cellcolor{blue!20}1 & \cellcolor{blue!20}1 & \cellcolor{blue!20}1 & \cellcolor{blue!20}1 & \cellcolor{blue!20}1 & 0 \\
        \hline
        0 & 0 & 0 & 0 & 0 & 0 & 0 & \cellcolor{blue!20}1 & \cellcolor{blue!20}1 & \cellcolor{blue!20}1 & 0 \\
        \hline
        0 & 0 & 0 & 0 & 0 & 0 & 0 & 0 & 0 & 0 & 0 \\
        \hline
        \end{tabular}
    \end{center}

    \begin{center}
        \begin{tabular}{|c|c|c|c|c|c|c|c|c|c|c|}
        \multicolumn{11}{c}{\textbf{Erosion}} \\
        \hline
        0 & 0 & 0 & 0 & 0 & 0 & 0 & 0 & 0 & 0 & 0 \\
        \hline
        0 & 0 & 0 & 0 & 0 & 0 & 0 & 0 & 0 & 0 & 0 \\
        \hline
        0 & 0 & 0 & 0 & \cellcolor{blue!20}1 & \cellcolor{blue!20}1 & \cellcolor{blue!20}1 & 0 & 0 & 0 & 0 \\
        \hline
        0 & 0 & 0 & 0 & 0 & 0 & \cellcolor{blue!20}1 & 0 & 0 & 0 & 0 \\
        \hline
        0 & 0 & 0 & 0 & 0 & 0 & 0 & 0 & 0 & 0 & 0 \\
        \hline
        0 & 0 & 0 & 0 & 0 & 0 & 0 & 0 & 0 & 0 & 0 \\
        \hline
        0 & 0 & 0 & 0 & 0 & 0 & 0 & 0 & 0 & 0 & 0 \\
        \hline
        0 & 0 & 0 & 0 & 0 & 0 & 0 & 0 & 0 & 0 & 0 \\
        \hline
        0 & 0 & 0 & 0 & 0 & 0 & 0 & 0 & 0 & 0 & 0 \\
        \hline
        0 & 0 & 0 & 0 & 0 & 0 & 0 & 0 & 0 & 0 & 0 \\
        \hline
        0 & 0 & 0 & 0 & 0 & 0 & 0 & 0 & 0 & 0 & 0 \\
        \hline
        \end{tabular}
    \end{center}

    \begin{center}
        \begin{tabular}{|c|c|c|c|c|c|c|c|c|c|c|}
        \multicolumn{11}{c}{\textbf{Opening}} \\
        \hline
        0 & 0 & 0 & 0 & 0 & 0 & 0 & 0 & 0 & 0 & 0 \\
        \hline
        0 & 0 & 0 & 0 & 0 & 0 & 0 & 0 & 0 & 0 & 0 \\
        \hline
        0 & 0 & 0 & \cellcolor{blue!20}1 & \cellcolor{blue!20}1 & \cellcolor{blue!20}1 & \cellcolor{blue!20}1 & \cellcolor{blue!20}1 & 0 & 0 & 0 \\
        \hline
        0 & 0 & 0 & \cellcolor{blue!20}1 & \cellcolor{blue!20}1 & \cellcolor{blue!20}1 & \cellcolor{blue!20}1 & \cellcolor{blue!20}1 & 0 & 0 & 0 \\
        \hline
        0 & 0 & 0 & 0 & 0 & \cellcolor{blue!20}1 & \cellcolor{blue!20}1 & \cellcolor{blue!20}1 & 0 & 0 & 0 \\
        \hline
        0 & 0 & 0 & 0 & 0 & 0 & 0 & 0 & 0 & 0 & 0 \\
        \hline
        0 & 0 & 0 & 0 & 0 & 0 & 0 & 0 & 0 & 0 & 0 \\
        \hline
        0 & 0 & 0 & 0 & 0 & 0 & 0 & 0 & 0 & 0 & 0 \\
        \hline
        0 & 0 & 0 & 0 & 0 & 0 & 0 & 0 & 0 & 0 & 0 \\
        \hline
        0 & 0 & 0 & 0 & 0 & 0 & 0 & 0 & 0 & 0 & 0 \\
        \hline
        0 & 0 & 0 & 0 & 0 & 0 & 0 & 0 & 0 & 0 & 0 \\
        \hline
        \end{tabular}
    \end{center}

    \begin{center}
        \begin{tabular}{|c|c|c|c|c|c|c|c|c|c|c|}
        \multicolumn{11}{c}{\textbf{Closing}} \\
        \hline
        0 & 0 & 0 & 0 & 0 & 0 & 0 & 0 & 0 & 0 & 0 \\
        \hline
        0 & 0 & 0 & 0 & 0 & 0 & 0 & 0 & 0 & 0 & 0 \\
        \hline
        0 & 0 & \cellcolor{blue!20}1 & \cellcolor{blue!20}1 & \cellcolor{blue!20}1 & \cellcolor{blue!20}1 & \cellcolor{blue!20}1 & \cellcolor{blue!20}1 & 0 & 0 & 0 \\
        \hline
        0 & 0 & \cellcolor{blue!20}1 & \cellcolor{blue!20}1 & \cellcolor{blue!20}1 & \cellcolor{blue!20}1 & \cellcolor{blue!20}1 & \cellcolor{blue!20}1 & 0 & 0 & 0 \\
        \hline
        0 & 0 & \cellcolor{blue!20}1 & \cellcolor{blue!20}1 & \cellcolor{blue!20}1 & \cellcolor{blue!20}1 & \cellcolor{blue!20}1 & \cellcolor{blue!20}1 & 0 & 0 & 0 \\
        \hline
        0 & 0 & \cellcolor{blue!20}1 & 0 & 0 & 0 & \cellcolor{blue!20}1 & \cellcolor{blue!20}1 & 0 & 0 & 0 \\
        \hline
        0 & 0 & \cellcolor{blue!20}1 & 0 & 0 & 0 & \cellcolor{blue!20}1 & \cellcolor{blue!20}1 & 0 & 0 & 0 \\
        \hline
        0 & 0 & 0 & \cellcolor{blue!20}1 & \cellcolor{blue!20}1 & \cellcolor{blue!20}1 & \cellcolor{blue!20}1 & \cellcolor{blue!20}1 & 0 & 0 & 0 \\
        \hline
        0 & 0 & 0 & 0 & 0 & 0 & 0 & 0 & \cellcolor{blue!20}1 & 0 & 0 \\
        \hline
        0 & 0 & 0 & 0 & 0 & 0 & 0 & 0 & 0 & 0 & 0 \\
        \hline
        0 & 0 & 0 & 0 & 0 & 0 & 0 & 0 & 0 & 0 & 0 \\
        \hline
        \end{tabular}
    \end{center}

    % Problem 2
    \item
    First, we calculate the closing using the following structuring element (with origin in the center):

    \begin{center}
        \textbf{Structuring Element (SE)} \\
        \begin{tabular}{|c|c|c|}
        \hline
        1 & 1 & 1 \\
        \hline
        1 & 1 & 1 \\
        \hline
        1 & 1 & 1 \\
        \hline
        \end{tabular}
    \end{center}

    This fills in all the gaps that a 3x3 square cannot fit into and results in this image:

    \begin{center}
        \begin{tabular}{|c|c|c|c|c|c|c|c|c|c|c|}
        \multicolumn{11}{c}{\textbf{Closing of I using SE (K)}} \\
        \hline
        0 & 0 & 0 & 0 & 0 & 0 & 0 & 0 & 0 & 0 & 0 \\
        \hline
        0 & 0 & 0 & 0 & 0 & 0 & 0 & 0 & 0 & 0 & 0 \\
        \hline
        0 & 0 & \cellcolor{blue!20}1 & \cellcolor{blue!20}1 & \cellcolor{blue!20}1 & \cellcolor{blue!20}1 & \cellcolor{blue!20}1 & \cellcolor{blue!20}1 & 0 & 0 & 0 \\
        \hline
        0 & 0 & \cellcolor{blue!20}1 & \cellcolor{blue!20}1 & \cellcolor{blue!20}1 & \cellcolor{blue!20}1 & \cellcolor{blue!20}1 & \cellcolor{blue!20}1 & 0 & 0 & 0 \\
        \hline
        0 & 0 & \cellcolor{blue!20}1 & \cellcolor{blue!20}1 & \cellcolor{blue!20}1 & \cellcolor{blue!20}1 & \cellcolor{blue!20}1 & \cellcolor{blue!20}1 & 0 & 0 & 0 \\
        \hline
        0 & 0 & \cellcolor{blue!20}1 & \cellcolor{blue!20}1 & \cellcolor{blue!20}1 & \cellcolor{blue!20}1 & \cellcolor{blue!20}1 & \cellcolor{blue!20}1 & 0 & 0 & 0 \\
        \hline
        0 & 0 & \cellcolor{blue!20}1 & \cellcolor{blue!20}1 & \cellcolor{blue!20}1 & \cellcolor{blue!20}1 & \cellcolor{blue!20}1 & \cellcolor{blue!20}1 & 0 & 0 & 0 \\
        \hline
        0 & 0 & 0 & \cellcolor{blue!20}1 & \cellcolor{blue!20}1 & \cellcolor{blue!20}1 & \cellcolor{blue!20}1 & \cellcolor{blue!20}1 & 0 & 0 & 0 \\
        \hline
        0 & 0 & 0 & 0 & 0 & 0 & 0 & 0 & \cellcolor{blue!20}1 & 0 & 0 \\
        \hline
        0 & 0 & 0 & 0 & 0 & 0 & 0 & 0 & 0 & 0 & 0 \\
        \hline
        0 & 0 & 0 & 0 & 0 & 0 & 0 & 0 & 0 & 0 & 0 \\
        \hline
        \end{tabular}
    \end{center}

    Now we use the following structuring element (with origin in the center) to calculate the opening of image K:

    \begin{center}
        \textbf{Structuring Element (SE2)} \\
        \begin{tabular}{|c|c|c|}
        \hline
        1 & 1 & 0 \\
        \hline
        1 & 0 & 0 \\
        \hline
        0 & 0 & 0 \\
        \hline
        \end{tabular}
    \end{center}

    The result of this operation is the image J. The opening removes pixels from the image wherever the structuring element
    cannot fit inside the image. The right angle shape of SE2 easily fits everywhere except the bottom right corner where
    it cuts off the two bottom-rightmost two pixels.

    % Problem 3
    \item
    The optimal threshold will be at the intensity where the probability of assigning the pixel to object or background
    is equal:

    \[
        \frac{1}{2}p_1(T)=\frac{1}{2}p_2(T)
    \]
    \[
        p_1(T)=p_2(T)
    \]

    If the PDFs are equal anywhere, it will be within the overlap at \(4 \leq x \leq 6\). Therefore, we
    substitute \(p_1(T)\) and \(p_2(T)\) as:

    \[
        -\frac{T}{18} + \frac{1}{3}=\frac{2}{25}T - \frac{8}{25}
    \]

    And simplify:
    \[
        \frac{6 - T}{18}=\frac{2T-8}{25}
    \]
    \[
        150 - 25T = 36T - 144
    \]
    \[
        294 = 61T
    \]
    \[
        T = \frac{294}{61}
    \]

    The optimal threshold is at \(T=4.82\).

    % Problem 4
    \item

    \begin{enumerate}
        % Problem 4a
        \item

        After one iteration the cluster assignment for each point is:
        \begin{center}
            \textbf{Cluster Assignment (1 iteration)} \\
            \begin{tabular}{|c|c|}
            \hline
            Point & Cluster Number \\
            \hline
            (-3, 4) & 1 \\
            \hline
            (-2, 1) & 3 \\
            \hline
            (-2, -4) & 2 \\
            \hline
            (0, -3) & 2 \\
            \hline
            (0, 0) & 3 \\
            \hline
            (2, 2) & 3 \\
            \hline
            (3, 1) & 3 \\
            \hline
            (4, 2) & 3 \\
            \hline
            \end{tabular}
        \end{center}

        And the cluster centers are at:
        \begin{center}
            \textbf{Cluster Centers (1 iteration)} \\
            \begin{tabular}{|c|c|}
            \hline
            Cluster Number & Center \\
            \hline
            1 & (-3, 4) \\
            \hline
            2 & (-1, -3.5) \\
            \hline
            3 & (1.4, 1.2) \\
            \hline
            \end{tabular}
        \end{center}

        After convergence, the cluster assignments will be:
        \begin{center}
            \textbf{Cluster Assignment (Convergence)} \\
            \begin{tabular}{|c|c|}
            \hline
            Point & Cluster Number \\
            \hline
            (-3, 4) & 1 \\
            \hline
            (-2, 1) & 1 \\
            \hline
            (-2, -4) & 2 \\
            \hline
            (0, -3) & 2 \\
            \hline
            (0, 0) & 3 \\
            \hline
            (2, 2) & 3 \\
            \hline
            (3, 1) & 3 \\
            \hline
            (4, 2) & 3 \\
            \hline
            \end{tabular}
        \end{center}

        And the cluster centers are at:
        \begin{center}
            \textbf{Cluster Centers (Convergence)} \\
            \begin{tabular}{|c|c|}
            \hline
            Cluster Number & Center \\
            \hline
            1 & (-2.5, 2.5) \\
            \hline
            2 & (-1, -3.5) \\
            \hline
            3 & (2.25, 1.25) \\
            \hline
            \end{tabular}
        \end{center}

        % Problem 4b
        \item
        After convergence, the cluster assignments will be:
        \begin{center}
            \textbf{Cluster Assignment (Convergence)} \\
            \begin{tabular}{|c|c|}
            \hline
            Point & Cluster Number \\
            \hline
            (-3, 4) & 1 \\
            \hline
            (-2, 1) & 1 \\
            \hline
            (-2, -4) & 2 \\
            \hline
            (0, -3) & 2 \\
            \hline
            (0, 0) & 2 \\
            \hline
            (2, 2) & 3 \\
            \hline
            (3, 1) & 3 \\
            \hline
            (4, 2) & 3 \\
            \hline
            \end{tabular}
        \end{center}

        And the cluster centers are at:
        \begin{center}
            \textbf{Cluster Centers (Convergence)} \\
            \begin{tabular}{|c|c|}
            \hline
            Cluster Number & Center \\
            \hline
            1 & (-2.5, 2.5) \\
            \hline
            2 & (-0.67, -2.33) \\
            \hline
            3 & (3, 1.67) \\
            \hline
            \end{tabular}
        \end{center}

        % Problem 4c
        \item
        The positions of the initial cluster centers can affect the number of iterations before convergence, the final cluster assignments,
        and the locations of the final cluster centers. The initial cluster centers from (a) took 3 iterations compared to (b)'s
        2 iterations. The point at (0,0) ends up in cluster 3 in (a), but cluster 2 in (b). This coincides with the different final cluster centers
        we see for clusters 2 and 3. If we assume that the initial centers for clusters 1 and 2 are fixed, then the final assignment for (0,0)
        appears to be based on whether the initial position for cluster 3 is closer to the pair of points with negative y-values or the triplet with
        positive x and y-values.

        % Problem 4d
        \item \ % Add a space to force a left alignment for the number.
        \begin{center}
            \textbf{Membership Assignments for (-3,4)} \\
            \begin{tabular}{|c|c|}
            \hline
            Cluster Number & Assignment \\
            \hline
            1 & 1 \\
            \hline
            2 & 0 \\
            \hline
            3 & 0 \\
            \hline
            \end{tabular}
        \end{center}

        \begin{center}
            \textbf{Membership Assignments for (0,0)} \\
            \begin{tabular}{|c|c|}
            \hline
            Cluster Number & Assignment \\
            \hline
            1 & .159 \\
            \hline
            2 & .442 \\
            \hline
            3 & .398 \\
            \hline
            \end{tabular}
        \end{center}

    \end{enumerate}
\end{enumerate}

\end{document}